\section{Gradient flow}
The gradient flow (GF) for gauge and fermion
fields~\cite{Luscher:2010iy,Luscher:2011bx,Luscher:2013cpa} provides
a different way to regulate short-distance singularities
of the twist-2 operators.
The connection with the physical renormalized matrix elements at vanishing flow time,
$t=0$, is obtained with a short-flow-time expansion (SFTX) after performing
the continuum limit of the hadronic matrix elements at a fixed
physical value of the flow time $t$~\cite{Luscher:2013vga}.

The flowed twist-2 operators, written in terms of flowed fermions, $\chi^r(x,t)$
and $\chibar^r(x,t)$, and gauge fields, $B_\mu(x,t)$, are given by
\be
O_n^{rs}(x,t) = \chibar^r(x,t) \gmuopen1 \lrDmu2 \cdots \lrDmuclose{n} \chi^s(x,t)\,,
\label{eq:flowed_t2}
\ee
where $\left\{\cdots\right\}$ denotes symmetrization over the included indices.
These operators are particularly advantageous because they renormalize
multiplicatively also on the lattice with a renormalization
factor that depends only on the fermion content of the operator,
\be
O_n^{rs}(t) = Z_n O_{n,B}^{rs}(t)\,, \quad Z_n = Z_\chi\,,
\ee
where $Z_\chi^{1/2}$ is the renormalization constant of the flowed
fermion fields~\cite{Luscher:2013cpa}.
For correlation functions containing flowed twist-2 operators,
the flow-time $t$ provides a regulator for short-distance singularities.
Beside the renormalization of the bare parameters of the theory,
the flowed fermion fields, and any local field at $t=0$,
correlation functions do not
require any additional renormalization.
For lattice QCD applications, it is convenient to define $Z_\chi$
in a scheme that is regularization independent, as then
the matching and the renormalization can be performed in the same scheme.
The choice of the specific scheme is not relevant for the general discussion,
but it becomes important when calculating the matching coefficients.
In the calculations of the matching coefficients described below,
we renormalize the twist-2 operators using the so-called ringed fields~\cite{Makino:2014taa},
defined by the SU($N_c$) gauge invariant and regularization independent condition
\be
\left \langle \rchibar_r(x,t) {\overset\leftrightarrow{\slashed{D}}} \rchi_r(x,t) \right\rangle
= - \frac{N_c}{(4 \pi)^2 t^2}\,.
\ee
Imposing this condition in dimensional regularization, with $D=4 - 2 \epsilon$,
leads to a finite renormalization between the ringed fields and
renormalized fields in the $\MS$ scheme
\bea
\chi_r(x,t) &=& \left(8 \pi t\right)^{\epsilon/2} \zeta_\chi^{1/2} \rchi_r(x,t) \nonumber \\
\chibar_r(x,t) &=& \left(8 \pi t\right)^{\epsilon/2} \zeta_\chi^{1/2} \rchibar_r(x,t)\,.
\eea
The 1-loop result for the finite renormalization is given by~\cite{Makino:2014taa}
\be
\zeta_\chi = 1 - \frac{\gbar^2(\mu)}{(4 \pi)^2}C_F
\left( 3 \log (8 \pi \mu^2 t) - \log (432)\right)\,,
\ee
where $\gbar$ is the strong coupling renormalized in the $\MS$ scheme
at the renormalization scale $\mu$, and $C_F = N_c^2-1/2N_c$.
This result has been extended to O($\gbar^4$) in
Refs.~\cite{Harlander:2018zpi,Artz:2019bpr}.
Results obtained in the $\MS$ scheme can be converted to the
$\MSbar$ scheme by replacing the renormalization scale $\mu$ with
$\mubar$ related by $\log \mu^2 = \log \mubar^2 + \gamma_E - \log 4 \pi$.

\subsection{O($a$) improvement}
The lattice QCD calculation of hadronic matrix elements of twist-2 operators
is affected by cutoff effects.

For nonperturbatively improved clover fermions,
the hadronic matrix elements of the standard
twist-2 operators (with $t=0$) are still affected
by O($a$) cutoff effects that could be removed by
adopting an O($a$) improved definition of the twist-2 operators.
Excluding the case $n=2$, where the
corresponding improvement coefficients have been determined
at 1-loop in perturbation theory~\cite{Capitani:2000xi},
the continuum limit is reached with O($a$)
discretization effects.

If we consider flowed twist-2 operators for a generic $n$,
the hadronic matrix elements
are affected only by O($a m$) cutoff effects,
where $m$ is the quark mass. Additional improvement
coefficients specific to the twist-2 operators are not needed.
In the discussion of the method's applications
below, we also consider
ratios of flowed matrix elements where both $Z_\chi$ and
the O($am$) cutoff effects are eliminated.
In summary, regardless of the lattice setup adopted,
the use of flowed twist-2 operators not only provides a significant simplification
in the renormalization pattern, but also
improves the parametric scaling in the lattice spacing
towards the continuum limit.
